\section{Problem Formulation}
Let \(R=\{r_1,r_2,\ldots,r_n\}\) be a batch of resumes submitted for a job description \(J\). The objective is to produce a ranked list \(L\) where each candidate receives a bounded match score \(S_i \in [0,100]\) and an explanation \(E_i\). The screening task is formulated as an auditable first-pass filtering problem, not as a fully autonomous hiring decision:
\begin{equation}
f(r_i,J,W,M) \rightarrow (S_i,E_i),
\end{equation}
where \(W\) denotes recruiter-configurable scoring weights and \(M\) denotes the selected contextual similarity mode. Each processed resume is represented as a structured record containing extracted text, detected technical skills, soft skills, inferred years of experience, contextual similarity, component scores, final score, and explanation. The job description is processed in the same way, but its extracted fields act as requirements. The design targets four requirements: format tolerance, bounded scoring, controllable computational cost, and human-readable interpretability. The framework deliberately avoids demographic inference and restricts ranking features to text-derived role relevance.
